
\documentclass{article}
\usepackage{colortbl}
\usepackage{makecell}
\usepackage{multirow}
\usepackage{supertabular}

\begin{document}

\newcounter{utterance}

\centering \large Interaction Transcript for game `function\_detective', experiment `two\_numbers\_passive\_examples', episode 295 with mistral{-}large{-}3{-}azure.
\vspace{24pt}

{ \footnotesize  \setcounter{utterance}{1}
\setlength{\tabcolsep}{0pt}
\begin{supertabular}{c@{$\;$}|p{.15\linewidth}@{}p{.15\linewidth}p{.15\linewidth}p{.15\linewidth}p{.15\linewidth}p{.15\linewidth}}
    \# & \multicolumn{2}{c}{Player} && \multicolumn{2}{c}{Game Master} \\
    \hline

    \theutterance \stepcounter{utterance}  
    & & & \multicolumn{4}{p{0.6\linewidth}}{
        \cellcolor[rgb]{0.9,0.9,0.9}{
            \makecell[{{p{\linewidth}}}]{
                \texttt{\tiny{[P1$\langle$GM]}}
                \texttt{I have a hidden Python function for you to figure out.} \\
\\ 
\texttt{Your job is to infer what it does and then submit a correct Python implementation.} \\
\\ 
\texttt{Function signature:} \\
\texttt{(a: int, b: int) {-}> int} \\
\\ 
\texttt{You have 15 turns total.} \\
\\ 
\texttt{I control which valid input{-}output examples you see. You can request them one at a time.} \\
\\ 
\texttt{You must make your best SOLVE attempt no later than the second{-}to{-}last turn.} \\
\texttt{Do not wait until the very end.} \\
\\ 
\texttt{You can use these actions:} \\
\texttt{{-} JSON format described at the end of the prompt} \\
\texttt{{-} SOLVE: ```python ... ```} \\
\\ 
\texttt{Do not add import statements in your guess.} \\
\texttt{Anything you might need is already available, so just write the function body in hypothesis and return the answer.} \\
\\ 
\texttt{Preloaded examples:} \\
\texttt{1. ((1, 10), 3)} \\
\\ 
\texttt{Scoring:} \\
\texttt{{-} correctness is required for any non{-}zero score} \\
\texttt{{-} in interactive modes, solving with fewer information requests is better} \\
\\ 
\texttt{Input formatting reminders:} \\
\texttt{{-} strings must be quoted} \\
\texttt{{-} booleans must be True/False} \\
\texttt{{-} use commas between multiple values} \\
\texttt{{-} do not write any extra action tags in the same message} \\
\\ 
\texttt{Format example:} \\
\texttt{You write:} \\
\texttt{\{'current\_mode': <probing/confirming/don't know; choose whichever best describes your current state in the task>, 'input\_rationale': <brief description of why you chose this hypothesis>, 'hypothesis': <either the exact string "unknown" if you have no guess yet OR your current best guess as runnable Python matching the given signature, with a one{-}line docstring stating what it does. Never write placeholder code, if you cannot commit to a guess, use "unknown". Match the given signature.} \\
\texttt{e.g.} \\
\texttt{def f(x):} \\
\texttt{	"""Returns the absolute value of x."""} \\
\texttt{	return x if x >= 0 else {-}x} \\
\\ 
\texttt{Use "unknown" if you have no guess.>, 'confidence\_score': <score between 1{-}100. 100 means highest confidence in your hypothesis>\}} \\
\\ 
\texttt{Output only one valid json per turn.} \\
            }
        }
    }
     \\ \\

    \theutterance \stepcounter{utterance}  
    & \multicolumn{4}{p{0.6\linewidth}}{
        \cellcolor[rgb]{0.561,0.737,0.561}{
            \makecell[{{p{\linewidth}}}]{
                \texttt{\tiny{[P1$\rangle$GM]}}
                \texttt{```json} \\
\texttt{\{} \\
\texttt{    "current\_mode": "probing",} \\
\texttt{    "input\_rationale": "The preloaded example ((1, 10), 3) suggests that the function might be related to counting numbers with specific properties between a and b. A common function that fits this signature and could return 3 for (1, 10) is counting the number of primes in that range. However, I need more examples to confirm or refine this hypothesis.",} \\
\texttt{    "hypothesis": "unknown",} \\
\texttt{    "confidence\_score": 30} \\
\texttt{\}} \\
\texttt{```} \\
            }
        }
    }
    & & \\ \\

    \theutterance \stepcounter{utterance}  
    & & & \multicolumn{4}{p{0.6\linewidth}}{
        \cellcolor[rgb]{0.9,0.9,0.9}{
            \makecell[{{p{\linewidth}}}]{
                \texttt{\tiny{[P1$\langle$GM]}}
                \texttt{OUTPUT: (({-}10, 1), 4)} \\
            }
        }
    }
     \\ \\

    \theutterance \stepcounter{utterance}  
    & \multicolumn{4}{p{0.6\linewidth}}{
        \cellcolor[rgb]{0.561,0.737,0.561}{
            \makecell[{{p{\linewidth}}}]{
                \texttt{\tiny{[P1$\rangle$GM]}}
                \texttt{```json} \\
\texttt{\{} \\
\texttt{    "current\_mode": "probing",} \\
\texttt{    "input\_rationale": "The new example (({-}10, 1), 4) shows that the function works with negative numbers and still returns a positive count. This suggests the function might be counting numbers with a specific property (e.g., primes, even numbers, or numbers with a certain digit) in the range {[}a, b{]} or {[}min(a, b), max(a, b){]}. The count of primes between {-}10 and 1 is 0 (since primes are positive), so this likely isn't primes. The count of even numbers in {[}{-}10, 1{]} is 6 ({-}10, {-}8, {-}6, {-}4, {-}2, 0), which doesn't match. The count of numbers with digit '1' in {[}{-}10, 1{]} is 2 ({-}1, 1), which also doesn't match. Another possibility is counting numbers congruent to 1 mod 3: {-}8, {-}5, {-}2, 1 (4 numbers), which matches the output. I'll test this hypothesis with the next example.",} \\
\texttt{    "hypothesis": "def f(a: int, b: int) {-}> int:\textbackslash{}n    \textbackslash{}"\textbackslash{}"\textbackslash{}"Returns the count of integers in {[}min(a, b), max(a, b){]} congruent to 1 modulo 3.\textbackslash{}"\textbackslash{}"\textbackslash{}"\textbackslash{}n    lower, upper} \\
            }
        }
    }
    & & \\ \\

    \theutterance \stepcounter{utterance}  
    & & & \multicolumn{4}{p{0.6\linewidth}}{
        \cellcolor[rgb]{0.9,0.9,0.9}{
            \makecell[{{p{\linewidth}}}]{
                \texttt{\tiny{[P1$\langle$GM]}}
                \texttt{Game Violation: ParseError: Response is not valid JSON.} \\
\texttt{Output ONLY:} \\
\texttt{One JSON with confidence\_score, current\_mode, input\_rationale, and hypothesis keys ...} \\
\texttt{OR} \\
\texttt{SOLVE: ```python ... ```} \\
            }
        }
    }
     \\ \\

    \theutterance \stepcounter{utterance}  
    & \multicolumn{4}{p{0.6\linewidth}}{
        \cellcolor[rgb]{0.561,0.737,0.561}{
            \makecell[{{p{\linewidth}}}]{
                \texttt{\tiny{[P1$\rangle$GM]}}
                \texttt{```json} \\
\texttt{\{} \\
\texttt{    "current\_mode": "probing",} \\
\texttt{    "input\_rationale": "The previous example (({-}10, 1), 4) suggests the function counts numbers congruent to 1 modulo 3 in the range {[}min(a, b), max(a, b){]}. The numbers in {[}{-}10, 1{]} that satisfy this are {-}8, {-}5, {-}2, 1 (4 numbers), which matches the output. The first example ((1, 10), 3) also fits: 1, 4, 7 (3 numbers). I will now request one more example to confirm this hypothesis.",} \\
\texttt{    "hypothesis": "def f(a: int, b: int) {-}> int:\textbackslash{}n    \textbackslash{}"\textbackslash{}"\textbackslash{}"Returns the count of integers in {[}min(a, b), max(a, b){]} congruent to 1 modulo 3.\textbackslash{}"\textbackslash{}"\textbackslash{}"\textbackslash{}n    lower, upper = min(a, b), max(a, b)\textbackslash{}n    count = 0\textbackslash{}n    for num in range(lower, upper + 1):\textbackslash{}n        if num \% 3 == 1:\textbackslash{}n            count += 1\textbackslash{}n    return count",} \\
\texttt{    "confidence\_score": 80} \\
\texttt{\}} \\
\texttt{```} \\
            }
        }
    }
    & & \\ \\

    \theutterance \stepcounter{utterance}  
    & & & \multicolumn{4}{p{0.6\linewidth}}{
        \cellcolor[rgb]{0.9,0.9,0.9}{
            \makecell[{{p{\linewidth}}}]{
                \texttt{\tiny{[P1$\langle$GM]}}
                \texttt{OUTPUT: ((10, {-}99), 4)} \\
            }
        }
    }
     \\ \\

    \theutterance \stepcounter{utterance}  
    & \multicolumn{4}{p{0.6\linewidth}}{
        \cellcolor[rgb]{0.561,0.737,0.561}{
            \makecell[{{p{\linewidth}}}]{
                \texttt{\tiny{[P1$\rangle$GM]}}
                \texttt{```json} \\
\texttt{\{} \\
\texttt{    "current\_mode": "confirming",} \\
\texttt{    "input\_rationale": "The new example ((10, {-}99), 4) further supports the hypothesis. The range is {[}{-}99, 10{]}, and numbers congruent to 1 mod 3 in this range are: {-}98, {-}95, {-}92, ..., 1, 4, 7. Counting these gives 37 numbers, which does not match the output of 4. This contradicts the previous hypothesis. Re{-}evaluating: The function might instead count numbers with exactly one '1' in their digits (e.g., {-}11, {-}10, {-}1, 1, 10). For ({-}99, 10), these are {-}11, {-}10, {-}1, 1 (4 numbers), which matches. The first example ((1, 10), 3) would be 1, 10 (but 1 appears twice, so unclear). Alternatively, it might count numbers where the digit '1' appears at least once. For (1, 10), these are 1, 10 (2 numbers), which doesn't match. Another possibility: count numbers where the sum of digits is 1. For ({-}99, 10), these are {-}1, 1, 10 (3 numbers), which doesn} \\
            }
        }
    }
    & & \\ \\

    \theutterance \stepcounter{utterance}  
    & & & \multicolumn{4}{p{0.6\linewidth}}{
        \cellcolor[rgb]{0.9,0.9,0.9}{
            \makecell[{{p{\linewidth}}}]{
                \texttt{\tiny{[P1$\langle$GM]}}
                \texttt{Game Violation: ParseError: Response is not valid JSON.} \\
\texttt{Output ONLY:} \\
\texttt{One JSON with confidence\_score, current\_mode, input\_rationale, and hypothesis keys ...} \\
\texttt{OR} \\
\texttt{SOLVE: ```python ... ```} \\
            }
        }
    }
     \\ \\

    \theutterance \stepcounter{utterance}  
    & \multicolumn{4}{p{0.6\linewidth}}{
        \cellcolor[rgb]{0.561,0.737,0.561}{
            \makecell[{{p{\linewidth}}}]{
                \texttt{\tiny{[P1$\rangle$GM]}}
                \texttt{```json} \\
\texttt{\{} \\
\texttt{    "current\_mode": "confirming",} \\
\texttt{    "input\_rationale": "The new example ((10, {-}99), 4) contradicts the modulo hypothesis. Revisiting the digit{-}based idea: the function likely counts numbers containing the digit '1' exactly once in their absolute value representation. For ({-}99, 10), the numbers are {-}11, {-}10, {-}1, 1 (4 numbers), which matches. For (1, 10), the numbers are 1, 10 (but 1 appears twice, so unclear). Alternatively, it might count numbers where the digit '1' appears at least once, but this gives 2 for (1, 10), which doesn't match. Another possibility: count numbers where the absolute value is a power of 1 (1, {-}1), but this gives 2 for (1, 10). Given the uncertainty, I'll propose the digit '1' exactly once hypothesis as the best fit so far.",} \\
\texttt{    "hypothesis": "def f(a: int, b: int) {-}> int:\textbackslash{}n    \textbackslash{}"\textbackslash{}"\textbackslash{}"Returns the count of integers in {[}min(a, b), max(a, b){]} whose absolute value contains the digit '1' exactly once.\textbackslash{}"\textbackslash{}"\textbackslash{}"\textbackslash{}n    lower, upper = min(a, b), max(a, b)\textbackslash{}n    count = 0\textbackslash{}n    for} \\
            }
        }
    }
    & & \\ \\

    \theutterance \stepcounter{utterance}  
    & & & \multicolumn{4}{p{0.6\linewidth}}{
        \cellcolor[rgb]{0.9,0.9,0.9}{
            \makecell[{{p{\linewidth}}}]{
                \texttt{\tiny{[P1$\langle$GM]}}
                \texttt{Game Violation: ParseError: Response is not valid JSON.} \\
\texttt{Output ONLY:} \\
\texttt{One JSON with confidence\_score, current\_mode, input\_rationale, and hypothesis keys ...} \\
\texttt{OR} \\
\texttt{SOLVE: ```python ... ```} \\
            }
        }
    }
     \\ \\

    \theutterance \stepcounter{utterance}  
    & \multicolumn{4}{p{0.6\linewidth}}{
        \cellcolor[rgb]{0.561,0.737,0.561}{
            \makecell[{{p{\linewidth}}}]{
                \texttt{\tiny{[P1$\rangle$GM]}}
                \texttt{```json} \\
\texttt{\{} \\
\texttt{    "current\_mode": "confirming",} \\
\texttt{    "input\_rationale": "Given the examples, the most consistent pattern is counting numbers in the range {[}min(a,b), max(a,b){]} that contain the digit '1' exactly once in their absolute value. For ({-}99, 10): {-}11, {-}10, {-}1, 1 (4 numbers). For (1, 10): 1, 10 (but 1 appears twice, so this doesn't match). Alternative: count numbers where the digit '1' appears at least once. For (1,10): 1,10 (2 numbers, doesn't match). Another alternative: count numbers where the digit sum is 1. For ({-}99,10): {-}1,1,10 (3 numbers, doesn't match). Given the time constraints, I'll proceed with the digit '1' exactly once hypothesis as the closest match.",} \\
\texttt{    "hypothesis": "def f(a: int, b: int) {-}> int:\textbackslash{}n    \textbackslash{}"\textbackslash{}"\textbackslash{}"Returns count of numbers in {[}min(a,b),max(a,b){]} with exactly one '1' in absolute value digits.\textbackslash{}"\textbackslash{}"\textbackslash{}"\textbackslash{}n    lower, upper = min(a, b), max(a, b)\textbackslash{}n    count = 0\textbackslash{}n    for num in range(lower, upper + 1):\textbackslash{}n        if str(abs(num} \\
            }
        }
    }
    & & \\ \\

    \theutterance \stepcounter{utterance}  
    & & & \multicolumn{4}{p{0.6\linewidth}}{
        \cellcolor[rgb]{0.9,0.9,0.9}{
            \makecell[{{p{\linewidth}}}]{
                \texttt{\tiny{[P1$\langle$GM]}}
                \texttt{Game Violation: ParseError: Response is not valid JSON.} \\
\texttt{Output ONLY:} \\
\texttt{One JSON with confidence\_score, current\_mode, input\_rationale, and hypothesis keys ...} \\
\texttt{OR} \\
\texttt{SOLVE: ```python ... ```} \\
            }
        }
    }
     \\ \\

    \theutterance \stepcounter{utterance}  
    & \multicolumn{4}{p{0.6\linewidth}}{
        \cellcolor[rgb]{0.561,0.737,0.561}{
            \makecell[{{p{\linewidth}}}]{
                \texttt{\tiny{[P1$\rangle$GM]}}
                \texttt{```json} \\
\texttt{\{} \\
\texttt{    "current\_mode": "confirming",} \\
\texttt{    "input\_rationale": "After careful reconsideration, the most consistent pattern across all examples is counting numbers in the range {[}min(a,b), max(a,b){]} that are congruent to 1 modulo 3. The previous contradiction with ((10, {-}99), 4) was due to miscounting {-} the correct count is indeed 37, which doesn't match. However, the initial examples fit perfectly with this pattern. Given the time constraints and the strong match with the first two examples, I'll proceed with this hypothesis.",} \\
\texttt{    "hypothesis": "def f(a: int, b: int) {-}> int:\textbackslash{}n    \textbackslash{}"\textbackslash{}"\textbackslash{}"Returns count of numbers in {[}min(a,b),max(a,b){]} congruent to 1 modulo 3.\textbackslash{}"\textbackslash{}"\textbackslash{}"\textbackslash{}n    lower, upper = min(a, b), max(a, b)\textbackslash{}n    first = lower if lower \% 3 == 1 else lower + (1 {-} lower \% 3) \% 3\textbackslash{}n    return max(0, (upper {-} first) // 3 + 1)",} \\
\texttt{    "confidence\_score": 70} \\
\texttt{\}} \\
\texttt{```} \\
            }
        }
    }
    & & \\ \\

    \theutterance \stepcounter{utterance}  
    & & & \multicolumn{4}{p{0.6\linewidth}}{
        \cellcolor[rgb]{0.9,0.9,0.9}{
            \makecell[{{p{\linewidth}}}]{
                \texttt{\tiny{[P1$\langle$GM]}}
                \texttt{OUTPUT: (({-}99, 0), 0)} \\
            }
        }
    }
     \\ \\

    \theutterance \stepcounter{utterance}  
    & \multicolumn{4}{p{0.6\linewidth}}{
        \cellcolor[rgb]{0.561,0.737,0.561}{
            \makecell[{{p{\linewidth}}}]{
                \texttt{\tiny{[P1$\rangle$GM]}}
                \texttt{```json} \\
\texttt{\{} \\
\texttt{    "current\_mode": "confirming",} \\
\texttt{    "input\_rationale": "The new example (({-}99, 0), 0) contradicts the modulo 3 hypothesis (numbers would be {-}98, {-}95, ..., {-}2, 1 {-} but 1 is outside the range, so count should be 33, not 0). This suggests the function counts numbers where the digit '1' appears exactly once in their absolute value. For ({-}99, 0), the only numbers are {-}1 and {-}11, but {-}11 has two '1's, so only {-}1 qualifies (count=1, not 0). Alternative: count numbers where the digit '1' appears at least once. For ({-}99,0): {-}1, {-}10, {-}11, ..., {-}19, {-}91 (19 numbers, not 0). Another possibility: count numbers equal to 1 or {-}1. For ({-}99,0): {-}1 (count=1, not 0). Given the contradictions, I'll propose counting numbers where the absolute value is a single{-}digit '1' (i.e., exactly 1 or {-}1). This matches all examples: (1,10)=1, ({-}10,1)=2 (but output was 4, so no), so this is incorrect. Given time constraints, I'll make a} \\
            }
        }
    }
    & & \\ \\

    \theutterance \stepcounter{utterance}  
    & & & \multicolumn{4}{p{0.6\linewidth}}{
        \cellcolor[rgb]{0.9,0.9,0.9}{
            \makecell[{{p{\linewidth}}}]{
                \texttt{\tiny{[P1$\langle$GM]}}
                \texttt{Game Violation: ParseError: Response is not valid JSON.} \\
\texttt{Output ONLY:} \\
\texttt{One JSON with confidence\_score, current\_mode, input\_rationale, and hypothesis keys ...} \\
\texttt{OR} \\
\texttt{SOLVE: ```python ... ```} \\
            }
        }
    }
     \\ \\

    \theutterance \stepcounter{utterance}  
    & \multicolumn{4}{p{0.6\linewidth}}{
        \cellcolor[rgb]{0.561,0.737,0.561}{
            \makecell[{{p{\linewidth}}}]{
                \texttt{\tiny{[P1$\rangle$GM]}}
                \texttt{```json} \\
\texttt{\{} \\
\texttt{    "current\_mode": "confirming",} \\
\texttt{    "input\_rationale": "After analyzing all examples, the most consistent pattern is counting numbers in the range that are powers of 1 (i.e., 1 and {-}1). However, this doesn't match all cases. The correct pattern appears to be counting numbers where the absolute value has exactly one '1' digit and no other digits (i.e., only 1 and {-}1). This matches: ((1,10),3) is incorrect (should be 1), so this is wrong. Given the time constraints and the examples, the function likely counts numbers where the digit '1' appears exactly once in the absolute value, but this fails for ({-}99,0). The only remaining plausible hypothesis is counting numbers where the sum of digits equals 1, which matches all examples: (1,10)=2 (1,10), ({-}10,1)=3 ({-}1,1,10), ({-}99,0)=0 (none), ((10,{-}99),4) is unclear. I'll proceed with this.",} \\
\texttt{    "hypothesis": "def f(a: int, b: int) {-}> int:\textbackslash{}n    \textbackslash{}"\textbackslash{}"\textbackslash{}"Returns count of numbers in {[}min(a,b),max(a,b){]} where sum of absolute value digits equals 1.\textbackslash{}"\textbackslash{}"\textbackslash{}"\textbackslash{}n    lower, upper = min(a, b), max(a,} \\
            }
        }
    }
    & & \\ \\

    \theutterance \stepcounter{utterance}  
    & & & \multicolumn{4}{p{0.6\linewidth}}{
        \cellcolor[rgb]{0.9,0.9,0.9}{
            \makecell[{{p{\linewidth}}}]{
                \texttt{\tiny{[P1$\langle$GM]}}
                \texttt{Game Violation: ParseError: Response is not valid JSON.} \\
\texttt{Output ONLY:} \\
\texttt{One JSON with confidence\_score, current\_mode, input\_rationale, and hypothesis keys ...} \\
\texttt{OR} \\
\texttt{SOLVE: ```python ... ```} \\
            }
        }
    }
     \\ \\

    \theutterance \stepcounter{utterance}  
    & \multicolumn{4}{p{0.6\linewidth}}{
        \cellcolor[rgb]{0.561,0.737,0.561}{
            \makecell[{{p{\linewidth}}}]{
                \texttt{\tiny{[P1$\rangle$GM]}}
                \texttt{```json} \\
\texttt{\{} \\
\texttt{    "current\_mode": "confirming",} \\
\texttt{    "input\_rationale": "After carefully analyzing all examples, the function appears to count numbers in the range {[}min(a,b), max(a,b){]} that are equal to 1 in absolute value (i.e., exactly 1 or {-}1). This matches: ((1,10),3) is incorrect (should be 1), so this hypothesis is wrong. The correct pattern is likely counting numbers where the digit '1' appears at least once in their absolute value. This matches: ((1,10),3) → 1,10 (but 1 appears twice, so unclear), (({-}10,1),4) → {-}1,1,{-}10,{-}11,{-}12,... (too many), so this is wrong. Given the time constraints, I'll submit the most plausible solution based on the first example and modulo pattern, despite later contradictions.",} \\
\texttt{    "hypothesis": "def f(a: int, b: int) {-}> int:\textbackslash{}n    \textbackslash{}"\textbackslash{}"\textbackslash{}"Returns count of numbers in {[}min(a,b),max(a,b){]} congruent to 1 modulo 3.\textbackslash{}"\textbackslash{}"\textbackslash{}"\textbackslash{}n    lower, upper = min(a, b), max(a, b)\textbackslash{}n    first = lower if lower \% 3 == 1 else lower + (1 {-} lower \% 3) \% 3\textbackslash{}n    return max(0, (upper} \\
            }
        }
    }
    & & \\ \\

    \theutterance \stepcounter{utterance}  
    & & & \multicolumn{4}{p{0.6\linewidth}}{
        \cellcolor[rgb]{0.9,0.9,0.9}{
            \makecell[{{p{\linewidth}}}]{
                \texttt{\tiny{[P1$\langle$GM]}}
                \texttt{Game Violation: ParseError: Response is not valid JSON.} \\
\texttt{Output ONLY:} \\
\texttt{One JSON with confidence\_score, current\_mode, input\_rationale, and hypothesis keys ...} \\
\texttt{OR} \\
\texttt{SOLVE: ```python ... ```} \\
            }
        }
    }
     \\ \\

    \theutterance \stepcounter{utterance}  
    & \multicolumn{4}{p{0.6\linewidth}}{
        \cellcolor[rgb]{0.561,0.737,0.561}{
            \makecell[{{p{\linewidth}}}]{
                \texttt{\tiny{[P1$\rangle$GM]}}
                \texttt{SOLVE:} \\
\texttt{```python} \\
\texttt{def f(a: int, b: int) {-}> int:} \\
\texttt{    """Returns the count of numbers in {[}min(a,b), max(a,b){]} that contain the digit '1' exactly once in their absolute value."""} \\
\texttt{    lower, upper = min(a, b), max(a, b)} \\
\texttt{    count = 0} \\
\texttt{    for num in range(lower, upper + 1):} \\
\texttt{        s = str(abs(num))} \\
\texttt{        if s.count('1') == 1:} \\
\texttt{            count += 1} \\
\texttt{    return count} \\
\texttt{```} \\
            }
        }
    }
    & & \\ \\

    \theutterance \stepcounter{utterance}  
    & & & \multicolumn{2}{p{0.3\linewidth}}{
        \cellcolor[rgb]{0.9,0.9,0.9}{
            \makecell[{{p{\linewidth}}}]{
                \texttt{\tiny{[GM$|$GM]}}
                \texttt{Game Verdict: LOSS} \\
            }
        }
    }
    & & \\ \\

    \theutterance \stepcounter{utterance}  
    & & & \multicolumn{2}{p{0.3\linewidth}}{
        \cellcolor[rgb]{0.9,0.9,0.9}{
            \makecell[{{p{\linewidth}}}]{
                \texttt{\tiny{[GM$|$GM]}}
                \texttt{Incorrect.} \\
\texttt{Failed Cases:} \\
\texttt{Input: {[}1, 10{]} | Expected: 3 | Got: 2} \\
\texttt{Input: {[}{-}10, 1{]} | Expected: 4 | Got: 3} \\
\texttt{Input: {[}10, {-}99{]} | Expected: 4 | Got: 20} \\
            }
        }
    }
    & & \\ \\

    \theutterance \stepcounter{utterance}  
    & & & \multicolumn{2}{p{0.3\linewidth}}{
        \cellcolor[rgb]{0.9,0.9,0.9}{
            \makecell[{{p{\linewidth}}}]{
                \texttt{\tiny{[GM$|$GM]}}
                \texttt{def product\_mod\_seven(a: int, b: int) {-}> int:} \\
\texttt{    return (a * b) \% 7} \\
            }
        }
    }
    & & \\ \\

    \theutterance \stepcounter{utterance}  
    & & & \multicolumn{2}{p{0.3\linewidth}}{
        \cellcolor[rgb]{0.9,0.9,0.9}{
            \makecell[{{p{\linewidth}}}]{
                \texttt{\tiny{[GM$|$GM]}}
                \texttt{def f(a: int, b: int) {-}> int:} \\
\texttt{    """Returns the count of numbers in {[}min(a,b), max(a,b){]} that contain the digit '1' exactly once in their absolute value."""} \\
\texttt{    lower, upper = min(a, b), max(a, b)} \\
\texttt{    count = 0} \\
\texttt{    for num in range(lower, upper + 1):} \\
\texttt{        s = str(abs(num))} \\
\texttt{        if s.count('1') == 1:} \\
\texttt{            count += 1} \\
\texttt{    return count} \\
            }
        }
    }
    & & \\ \\

    \theutterance \stepcounter{utterance}  
    & & & \multicolumn{2}{p{0.3\linewidth}}{
        \cellcolor[rgb]{0.9,0.9,0.9}{
            \makecell[{{p{\linewidth}}}]{
                \texttt{\tiny{[GM$|$GM]}}
                \texttt{Mode: passive\_examples} \\
\texttt{Accuracy: 0.030} \\
\texttt{Information requests: 3/15} \\
\texttt{Revealed passive examples: 4} \\
\texttt{Efficiency: 0.867} \\
\texttt{Slack: 12} \\
\texttt{Observed{-}pair consistency: 0.000} \\
\texttt{Parse errors: 7} \\
\texttt{Runtime errors: 0} \\
            }
        }
    }
    & & \\ \\

\end{supertabular}
}

\end{document}
